\section*{Bab \ref{ch:b1:active-margins} --- Margin aktif}

\begin{solution}{exo:b1:active-margins:1}
(a) $2\pi G\,\Delta\rho\,h = 2\pi\times\num{6.674e-11}\times1700\times3000 =
\qty{2.14e-3}{m.s^{-2}}$: gravitasi akan lebih rendah sekitar
\qty{214}{\milli\gal}. (b) Untuk \qty{3.5}{km},
$2\pi\times\num{6.674e-11}\times1700\times3500 = \qty{2.50e-3}{m.s^{-2}}$,
\qty{250}{\milli\gal}: besaran rendahan yang teramati, sehingga palung jauh dari
kesetimbangan isostatik. (c) Rumus slab mengandaikan batuan yang hilang meluas tanpa
batas di setiap sisi; palung yang lebarnya hanya puluhan kilometer menghilangkan massa
jauh lebih sedikit di sekitar titik pengukuran, sehingga airnya saja memberi kurang dari
\qty{250}{\milli\gal}. Sisa defisitnya terletak lebih dalam: seluruh lempeng dibengkokkan
ke bawah di bawah palung, dan \omterm{def:b1:layered-earth:layers}{keraknya}, yang lebih ringan daripada \omterm{def:b1:layered-earth:layers}{mantel}, menempati
tempat \omterm{def:b1:layered-earth:layers}{mantel} di dasarnya.
\end{solution}

\begin{solution}{exo:b1:active-margins:2}
(a) $\qty{82}{mm.yr^{-1}} = \qty{82}{km.Ma^{-1}}$, sehingga \qty{820}{km} dalam
\qty{10}{Ma}. (b) $820\sin\ang{45} = \qty{580}{km}$. (c) Lintasan ke
\qty{600}{km} adalah $600/\sin\ang{45} = \qty{849}{km}$, ditempuh dalam
$849/82 = \qty{10.3}{Ma}$.
\end{solution}

\begin{solution}{exo:b1:active-margins:3}
$g = 110/\tan\delta$: \qty{191}{km} untuk \ang{30}, \qty{110}{km} untuk \ang{45},
\qty{64}{km} untuk \ang{60}. Busur di atas slab yang landai miring paling jauh dari
palungnya.
\end{solution}

\begin{solution}{exo:b1:active-margins:4}
(a) Konvergensi: $\qty{2.7e6}{m^2.yr^{-1}}/\qty{5.1e7}{m} =
\qty{0.053}{m.yr^{-1}}$, sekitar \qty{53}{mm.yr^{-1}}. Penyebaran:
$\qty{3.0e6}{m^2.yr^{-1}}/\qty{6.7e7}{m} = \qty{0.045}{m.yr^{-1}}$, sekitar
\qty{45}{mm.yr^{-1}}. Dengan panjang total yang lebih pendek, batas subduksi
menghabiskan dasar samudra lebih cepat daripada punggungan membuatnya, namun
menghabiskan lebih sedikit seluruhnya. (b) Sekitar \qty{0.3}{km^2} sisanya diserap
sepanjang batas konvergen lain dalam model, seperti tumbukan benua dan zona konvergen
di dalam samudra, tempat lempeng-lempeng memendek alih-alih menunjam; dan punggungan
bukan satu-satunya batas divergen: \omterm{def:b1:rifts-ridges-margins:rift}{celah benua} dan beberapa transform juga membuka
sedikit.
\end{solution}

\begin{solution}{exo:b1:active-margins:5}
(a) $\bar x = 350$, $\bar z = 290$; simpangan $x$ adalah $-150$, $-50$,
$50$, $150$ dan simpangan $z$ adalah $-140$, $-40$, $40$, $140$, sehingga
$S_{xx} = 2(150^2 + 50^2) = \num{50000}$ dan
$S_{xz} = 2(150\times140 + 50\times40) = \num{46000}$. Maka $s = 0.92$,
$a = 290 - 0.92\times350 = \qty{-32}{km}$, dan \omterm{def:b1:geological-maps:attitude}{kemiringannya}
$\arctan 0.92 = \ang{42.6}$. (b) $z = 0$ pada $x = 32/0.92 = \qty{34.8}{km}$
dari palung. (c) Melalui titik ujung, \omterm{def:b1:geological-maps:attitude}{kemiringannya} adalah $280/300 = 0.933$, sebuah
\omterm{def:b1:geological-maps:attitude}{kemiringan} \ang{43.0}: dekat, tetapi hanya memakai dua dari empat gempa bumi.
\end{solution}

\begin{solution}{exo:b1:active-margins:6}
(a) $50/\tan\ang{25} + 60/\tan\ang{45} = 107.2 + 60.0 = \qty{167.2}{km}$.
(b) Slab \ang{45} sepanjang jalannya memberi \qty{110}{km}: segmen dangkal yang landai
memperlebar celah sebesar \qty{57}{km}. (c) Segmen dalam adalah
$z = 50 + (x - 107.2)$; ia mencapai $z = 0$ pada $x = \qty{57.2}{km}$, di sisi daratan
dari palung, seperti garis Kuril pada \qty{53}{km}: segmen dangkal yang landai
menjelaskan pergeseran Kuril. (d) $107.2 + 60/\tan\ang{60} = 107.2 + 34.6 =
\qty{141.9}{km}$.
\end{solution}

\begin{solution}{exo:b1:active-margins:7}
(a) Jari-jari adalah setengah \omterm{def:b1:geological-maps:attitude}{kemiringan}, dan $r = R\,\theta$ dengan $\theta$ dalam radian:
\ang{22.5} dan $6371\times0.3927 = \qty{2502}{km}$ untuk \ang{45}; \ang{15} dan
\qty{1668}{km} untuk \ang{30}; \ang{35} dan \qty{3892}{km} untuk \ang{70}.
(b) Busur di atas slab yang landai. (c) Prediksinya adalah
$44.9/2 = \ang{22.45}$; kecocokannya, \ang{21.9}, adalah
$0.55/22.45 = \qty{2.4}{\percent}$ lebih kecil.
\end{solution}

\begin{solution}{exo:b1:active-margins:8}
(a) $q/k$: $0.035/2.5 = \qty{0.014}{K.m^{-1}} = \qty{14}{K.km^{-1}}$ di bawah
\omterm{def:b1:active-margins:active-margin}{forearc}, $0.080/2.5 = \qty{32}{K.km^{-1}}$ di bawah \omterm{def:b1:active-margins:active-margin}{cekungan busur-belakang}.
(b) $10 + 14\times30 = \qty{430}{\celsius}$ dan
$10 + 32\times30 = \qty{970}{\celsius}$. (c) Sebagian \omterm{def:b1:heat-budget:heat-flow}{aliran panas} permukaan dihasilkan
oleh radioaktivitas di dalam \omterm{def:b1:layered-earth:layers}{kerak} itu sendiri, sehingga lebih sedikit panas menyeberangi
tingkat-tingkat yang lebih dalam dan gradiennya berkurang dengan kedalaman: ekstrapolasi
linear dari gradien permukaan melebihkan suhu pada kedalaman.
\end{solution}

\begin{solution}{exo:b1:active-margins:9}
(a) \qty{32}{mm.yr^{-1}} adalah \qty{32}{km.Ma^{-1}}; setengah lajunya adalah
\qty{16}{mm.yr^{-1}}. (b) $100/32 = \qty{3.1}{Ma}$. (c) Laju itu tidak harus tetap;
sebagian lebar mungkin berasal dari peregangan dan pemekaran busur sebelum penyebaran
dimulai; dan sumbunya mungkin telah melompat atau menyebar secara asimetris.
\end{solution}

\begin{solution}{exo:b1:active-margins:10}
(a) Lintasannya adalah $600/\sin\ang{45} = \qty{849}{km}$, ditempuh pada
\qty{82}{km.Ma^{-1}} dalam \qty{10.3}{Ma}, yaitu \qty{3.27e14}{s}.
(b) $\sqrt{\num{e-6}\times\num{3.27e14}} = \qty{1.8e4}{m}$, sekitar
\qty{18}{km}. (c) $(\qty{5e4}{m})^2/\qty{e-6}{m^2.s^{-1}} = \qty{2.5e15}{s}$,
sekitar \qty{79}{Ma}. (d) Panas hanya mencapai \qty{18}{km} terluar dari setiap sisi
slab: intinya tiba di \qty{600}{km} hampir sedingin ketika ia meninggalkan palung, lebih
dingin dan lebih kuat daripada \omterm{def:b1:layered-earth:layers}{mantel} di sekitarnya, itulah sebabnya ia masih dapat
pecah dalam gempa bumi pada kedalaman itu.
\end{solution}

\begin{solution}{exo:b1:active-margins:11}
(a) Pada jarak $L$ di belakang ujung, permukaan berada $L\tan\alpha$ di atas tingkat
ujung dan dasarnya terletak $L\tan\beta$ di bawahnya: punggung prisma berketinggian
$L(\tan\alpha + \tan\beta)$, dan segitiganya memiliki luas
$\frac12 L\times L(\tan\alpha + \tan\beta)$. (b) $\tan\ang{4} + \tan\ang{6} =
0.175$. Prisma menerima $\qty{1}{km}\times\qty{40}{km.Ma^{-1}} =
\qty{40}{km^2.Ma^{-1}}$ penampang, sehingga $A = 40t$ dan
$L = \sqrt{2A/0.175}$: \qty{21.4}{km} setelah \qty{1}{Ma}, \qty{42.8}{km}
setelah \qty{4}{Ma}, \qty{67.6}{km} setelah \qty{10}{Ma}. $L$ tumbuh sebagai
$\sqrt t$: empat kali lebih lama memberi dua kali lebar. (c) $A$ dibagi dua dan $L$
dibagi $\sqrt2$: \qty{15.1}{km}, \qty{30.2}{km} dan \qty{47.8}{km}.
\end{solution}

\begin{solution}{exo:b1:active-margins:12}
(a) $\partial F/\partial a = -2\sum(z_i - a - s x_i) = 0$ memberi
$n\bar z = na + sn\bar x$, yaitu $\bar z = a + s\bar x$. (b) Setiap $z_i$ dan
$\bar z$ bertambah \qty{10}{km}, sehingga simpangan $z_i - \bar z$ dan $S_{xz}$
tidak berubah: \omterm{def:b1:geological-maps:attitude}{kemiringannya} sama, dan $a = \bar z - s\bar x$ bertambah
\qty{10}{km}. (c) Simpangan $x_i - \bar x$ tidak berubah, sehingga \omterm{def:b1:geological-maps:attitude}{kemiringannya} sama;
$\bar x$ berkurang 20, sehingga $a$ bertambah $20s$. (d) Tidak keduanya: \omterm{def:b1:geological-maps:attitude}{kemiringan}
hanya bergantung pada gradien. Kesalahan yang dimiliki bersama oleh semua titik
memindahkan garis tanpa memutarnya.
\end{solution}

\begin{solution}{pb:b1:active-margins:1}
\textbf{1.} Kedelapan titik berada dekat garis lurus yang miring sekitar
\ang{45}, dari \qty{96}{km} dalamnya pada \qty{155}{km} dari palung hingga
\qty{573}{km} dalamnya pada \qty{609}{km}.

\textbf{2.} $\sum x_i = 3152$, $\sum z_i = 2751$, $\sum x_i^2 = \num{1442910}$,
$\sum x_i z_i = \num{1290548}$; $\bar x = \qty{394.0}{km}$,
$\bar z = \qty{343.875}{km}$.

\textbf{3.} $S_{xx} = 1442910 - 8\times394^2 = \num{201022}$ dan
$S_{xz} = 1290548 - 8\times394\times343.875 = \num{206654}$. Maka
$s = 206654/201022 = 1.028$ dan $a = 343.875 - 1.02802\times394 =
\qty{-61.2}{km}$: $z = -61.2 + 1.028\,x$.

\textbf{4.} $\arctan1.028 = \ang{45.8}$.

\textbf{5.} Pada $x = 61.2/1.028 = \qty{59.5}{km}$ di sisi daratan dari palung. Slab
yang meninggalkan palung pada \omterm{def:b1:geological-maps:attitude}{kemiringan} ini akan mencapai permukaan di sana; tetapi ia
mencapainya di palung, sehingga ia miring lebih landai di dekat palung dan menjadi lebih
curam di bawahnya.

\textbf{6.} Simpangannya adalah $-2.1$, $2.0$, $-7.3$, $5.3$, $7.1$, $2.3$,
$-15.0$ dan $8.1$ \unit{km}. Peristiwa 7 terletak \qty{15}{km} di atas garis,
peristiwa 8 \qty{8}{km} di bawahnya; yang lain terletak dalam sekitar \qty{7}{km} darinya,
sebagian kecil dari ketebalan slab.

\textbf{7.} $\ang{320.9} - \ang{300} = \ang{20.9}$. Tegak lurus terhadap palung:
$87.4\cos\ang{20.9} = \qty{81.6}{mm.yr^{-1}}$; sejajar:
$87.4\sin\ang{20.9} = \qty{31.2}{mm.yr^{-1}}$. Konvergensinya miring.

\textbf{8.} \qty{81.6}{mm.yr^{-1}} adalah \qty{81.6}{km.Ma^{-1}}: \qty{816}{km} dalam
\qty{10}{Ma}.

\textbf{9.} Lintasannya adalah $632/\sin\ang{45} = \qty{894}{km}$, ditempuh dalam
$894/81.6 = \qty{11.0}{Ma}$.

\textbf{10.} $122 + 11 = \qty{133}{Ma}$.

\textbf{11.} $t = \qty{10.95}{Ma} = \qty{3.46e14}{s}$ dan
$\sqrt{\num{e-6}\times\num{3.46e14}} = \qty{1.86e4}{m}$, sekitar \qty{19}{km}
dari masing-masing sisi: jauh lebih sedikit daripada \qty{50}{km} ke tengah slab
setebal \qty{100}{km}, yang intinya masih dingin.

\textbf{12.} $q/k$: \qty{17.2}{K.km^{-1}} di bawah dasar samudra tua,
\qty{14.0}{K.km^{-1}} di bawah \omterm{def:b1:active-margins:active-margin}{forearc} dan \qty{32.0}{K.km^{-1}} di bawah
\omterm{def:b1:active-margins:active-margin}{cekungan busur-belakang}.

\textbf{13.} Di bawah \omterm{def:b1:active-margins:active-margin}{forearc} slab yang dingin hanya berada beberapa puluh kilometer ke
bawah; ia menyerap panas dari batuan di atasnya, sehingga lebih sedikit panas mencapai
permukaan.

\textbf{14.} $h = \qty{2.60e-3}{m.s^{-2}}/(2\pi\times\num{6.674e-11}\times1700)
= \qty{3647}{m}$, sekitar \qty{3.6}{km}: sebesar, atau sedikit lebih besar daripada,
seluruh kedalaman palung di bawah dasar di sekitarnya, sekitar \qty{3.5}{km}. Palung itu
kehilangan massa yang tidak dikompensasi oleh apa pun.

\textbf{15.} Selisih massa jenisnya adalah $3300 - 1000 = \qty{2300}{kg.m^{-3}}$,
dan $2\pi\times\num{6.674e-11}\times2300\times3500 = \qty{3.38e-3}{m.s^{-2}}$,
\qty{338}{\milli\gal}. Rumus slab menghitung massa yang hilang hingga jarak berapa pun;
palung yang sebenarnya hanya membengkok ke bawah pada pita yang lebarnya puluhan
kilometer, sehingga defisitnya lebih kecil, sekitar \qty{260}{\milli\gal}, di antara
kedua perkiraan slab.

\textbf{16.} Kekuatan lempeng yang membengkok, yang menahan palung ke bawah dan
tonjolan luar ke atas seperti lempeng elastik, dan tarikan slab padat yang tenggelam ke
dalam \omterm{def:b1:layered-earth:layers}{mantel}.

\textbf{17.} $-61.2 + 1.028\times164 = \qty{107.4}{km}$.

\textbf{18.} Ya: \qty{107}{km} berada di antara 72 dan \qty{173}{km}, dan hanya sekitar
\qty{2}{km} lebih dalam daripada rata-rata, \qty{105}{km}.

\textbf{19.} $x = (100 + 61.2)/1.028 = \qty{156.8}{km}$.

\textbf{20.} $\ang{45.8}/2 = \ang{22.9}$, dan
$6371\times22.9\times\pi/180 = \qty{2546}{km}$.

\textbf{21.} Kecocokannya, \ang{21.9}, adalah \ang{1.0} lebih kecil, sekitar
\qty{4}{\percent}: kulit bola yang dibengkokkan menggambarkan lengkungan busur Kuril
dengan baik.

\textbf{22.} $x = (105 + 61.2)/1.028 = 166.2/1.028 = \qty{161.7}{km}$.

\textbf{23.} $164 - 161.7 = \qty{2.3}{km}$ lebih dekat ke palung daripada garis depan
yang teramati, $2.3/164 = \qty{1.4}{\percent}$: gunung api Kuril berdiri di atas slab
yang sedikit lebih dalam daripada rata-rata, seperti yang ditemukan pada soal 18.

\textbf{24.} Hingga kilometer terdekat, garis slab Kuril memprediksi \omterm{def:b1:active-margins:volcanic-arc}{celah busur--palung}
sebesar \textbf{\qty[text-series-to-math=true]{162}{km}} untuk gunung api yang berdiri di
atas slab pada \qty{105}{km} (teramati: \qty{164}{km}).
\end{solution}
